% !PRINTED 185
\weberpage{185}

\begin{center}
\textbf{\S\,56. \gesperrt{Potenzreihen für die Weierstrassschen Funktionen}}\\
{\boldmath\textbf{$\wp(u), \sigma(u)$.}}
\end{center}

Die Funktion $\sigma(u)$, als Funktion von $u$ betrachtet, läßt sich,
wie schon bemerkt, in eine in der ganzen $u$-Ebene konvergierende
Reihe nach Potenzen von $u$ entwickeln, und nach \S\,35, (8), (12)
hat diese Entwickelung die folgende Form:

\weq{1}{\sigma(u) = u + \alpha u^3 + \alpha_1 u^5 + \alpha_2 u^7 + \alpha_3 u^9 + \cdots}

Die Funktion

\weq{2}{\wp(u) = -\frac{d^2\log\sigma(u)}{du^2} = \frac{\sigma'(u)^2 - \sigma\sigma''(u)}{\sigma(u)^2}}

läßt sich ebenfalls nach steigenden Potenzen von $u$ entwickeln,
aber nicht mehr in eine immer konvergente Reihe, sondern diese
Reihe hat einen Konvergenzkreis. Sie hat die Form:

\weq{3}{\wp(u) = \frac{1}{u^2} + a + a_1 u^2 + a_2 u^4 + a_3 u^6 + \cdots}

und man kann die Koeffizienten $a, a_1, a_2, a_3, \ldots$ der Reihe nach
aus der Differentialgleichung [\S\,46, (18)]

\weq{4}{\wp'(u)^2 = 4\wp(u)^3 - g_2\wp(u) - g_3}

als Funktionen von $g_2, g_3$ bestimmen.
Es ist zunächst
\[
\wp'(u) = -\frac{2}{u^3} + 2a_1 u + 4a_2 u^3 + 6a_3 u^5 + \cdots,
\]
\[
\wp'(u)^2 = \frac{4}{u^6} - \frac{8a_1}{u^2} - 16a_2 + (4a_1^2 - 24a_3)u^2 + \cdots,
\]
und da $\wp'(u)^2$ kein Glied mit $u^{-4}$ enthält, während in $\wp(u)^3$
das Glied $3au^{-4}$ vorkommt, so muß $a=0$ sein. Danach ergibt
sich:
\[
\begin{aligned}
\wp(u)^3 &= \frac{1}{u^6} + \frac{3a_1}{u^2} + 3a_2 + 3(a_3+a_1^2)u^2 + \cdots,\\
\wp(u) &= \frac{1}{u^2} \phantom{{}+ \frac{3a_1}{u^2} + 3a_2{}} + a_1 u^2 + \cdots
\end{aligned}
\]
Und wenn man dies in die Gleichung (4) einsetzt, so folgt

\weq{5}{a_1 = \frac{g_2}{20}, \quad a_2 = \frac{g_3}{28}, \quad a_3 = \frac{g_2^2}{1200}.}
